\scanpage{100}
% archive sheet number: 99
\deleted{Il y a un \unsure{s-gpe|ss-gpe|s-groupe} d'indice fini \inserted{$I'$} de \deleted{$\unsure{\pi_1(H)|\pi_1(U)|\pi_1(\mathcal U)}$}\inserted{$I$} qui opère sur $E$ via un pro-$\ell$-groupe. A fortiori,}

\deleted{\diagram{\begin{tikzcd}[ampersand replacement=\&]
 I_0\deleted{(\unsure{X,\xi})}\arrow[r]\arrow[d] \& I'_0\arrow[d]\\
 \deleted{\unsure{\pi}}I\deleted{(\unsure{X,\xi})}\arrow[r] \& I'
\end{tikzcd}}}

\textcircled{a} Un résultat \inserted{local} de Abhyankar \textcircled{$+$}

\textcircled{b} \uline{Version \inserted{semi-}globale du \deleted{\unsure{un}} résultat de Abhyankar}.

$X$ schéma loc. noeth. connexe

$\deleted{Y}D=\sum\deleted{Y_i}D_i$ diviseur \inserted{strictement} à croisements normaux
(composantes $Y_i$ \uline{régulières})

On suppose $Z=\bigcap_i\deleted{Y_i}D_i\ne\varnothing$, et que \deleted{\unsure{pour}} \inserted{le couple} $(X,Z)$ est ``hensélien'', i.e.
\unsure{p.r un|pour} tt $X'$ fini\inserted{ sur $X$, conn $\ne\varnothing$}, \deleted{\unsure{l'im. réciproque}} de $Z$ étant $Z'$, \deleted{$\pi_0(Z')\simeq\pi_0(X')$} \inserted{$Z'$ est connexe $\ne\varnothing$} \unsure{$\longleftrightarrow$}.

\deleted{Alors} Soit \deleted{$X$} $U=X-\deleted{\unsure{\operatorname{supp}}}Y$, $\xi$ un pt géom de $U$.

Alors a) $U$ connexe, et $\pi_1(U,\xi)_t\longrightarrow\pi_1(X,\xi)$ surjectif.

b) Pour tout $D_i$, le groupe d'inertie dans
$\pi_1(U,\xi)_t$ relatif à $\pi_1(U)$ est un sous-groupe distingué et
\deleted{$\pi_1(U)$ \unsure{qu'il}} commutatif $I_i$ de $Y_i$ \deleted{\unsure{on a}} et on a
un épimorphisme canonique
\[
 \deleted{\unsure{T}}\prod_{\substack{\ell\ \text{inversible}\\\text{sur }\deleted{\unsure{Z}}D_i\\(\text{au moins sur }Z)}}T_\ell(\Omega^*)\unsure{\longrightarrow|-}I_i .
\]
c) Les groupes $I_i$ commutent, et \deleted{leur produit} \inserted{le groupe engendré} est le noyau $I_\xi$ de $\pi_1(U,\xi)_t\longrightarrow\pi_1(X,\xi)$.

a) Est immédiat.

b) L'image inverse de $D_i$ dans un rev. \deleted{\unsure{étale}} \inserted{connexe} modérément
ramifié $X'$ de $X$ est normal (étude locale) et \uline{connexe},
\deleted{\unsure{donc $I_0$ est}} irréd, donc $I_i$ est invariant. On sait
bien qu'il est $\simeq$ quotient de $\prod_\ell T_\ell(\Omega^*)$, $\ell\ne$ car. gén. de $D_i$.
Il faut prouver que $I$ est premier aux car. résiduelles
de $D_i$. C'est une histoire locale sur $D_i$

c) Commutation des $I_i$ : \deleted{\unsure{Ok.}} Étude locale en un pt de
